\documentclass[10pt,a4paper]{amsart}
\usepackage{amsmath,amssymb,amsthm,graphicx}
% For commutative diagrams we use Xy-pic
\usepackage[all]{xy}

% We use 2cell for 2-commutative diagrams.
\xyoption{2cell}
% Theorem environments.
%
\theoremstyle{plain}
\newtheorem{theorem}[subsection]{Theorem}
\newtheorem{proposition}[subsection]{Proposition}
\newtheorem{lemma}[subsection]{Lemma}

\theoremstyle{definition}
\newtheorem{definition}[subsection]{Definition}
\newtheorem{example}[subsection]{Example}
\newtheorem{exercise}[subsection]{Exercise}
\newtheorem{situation}[subsection]{Situation}

\theoremstyle{remark}
\newtheorem{remark}[subsection]{Remark}
\newtheorem{remarks}[subsection]{Remarks}

\numberwithin{equation}{subsection}

% Macros
%
\def\lim{\mathop{\mathrm{lim}}\nolimits}
\def\colim{\mathop{\mathrm{colim}}\nolimits}
\def\Spec{\mathop{\mathrm{Spec}}}
\def\Hom{\mathop{\mathrm{Hom}}\nolimits}
\def\Ext{\mathop{\mathrm{Ext}}\nolimits}
\def\SheafHom{\mathop{\mathcal{H}\!\mathit{om}}\nolimits}
\def\SheafExt{\mathop{\mathcal{E}\!\mathit{xt}}\nolimits}
\def\Sch{\mathit{Sch}}
\def\Mor{\mathop{\mathrm{Mor}}\nolimits}
\def\Ob{\mathop{\mathrm{Ob}}\nolimits}
\def\Sh{\mathop{\mathit{Sh}}\nolimits}
\def\NL{\mathop{N\!L}\nolimits}
\def\CH{\mathop{\mathrm{CH}}\nolimits}
\def\proetale{{pro\text{-}\acute{e}tale}}
\def\etale{{\acute{e}tale}}
\def\QCoh{\mathit{QCoh}}
\def\Ker{\mathop{\mathrm{Ker}}}
\def\Im{\mathop{\mathrm{Im}}}
\def\Coker{\mathop{\mathrm{Coker}}}
\def\Coim{\mathop{\mathrm{Coim}}}

% Boxtimes
%
\DeclareMathSymbol{\boxtimes}{\mathbin}{AMSa}{"02}

%
% Macros for moduli stacks/spaces
%
\def\QCohstack{\mathcal{QC}\!\mathit{oh}}
\def\Cohstack{\mathcal{C}\!\mathit{oh}}
\def\Spacesstack{\mathcal{S}\!\mathit{paces}}
\def\Quotfunctor{\mathrm{Quot}}
\def\Hilbfunctor{\mathrm{Hilb}}
\def\Curvesstack{\mathcal{C}\!\mathit{urves}}
\def\Polarizedstack{\mathcal{P}\!\mathit{olarized}}
\def\Complexesstack{\mathcal{C}\!\mathit{omplexes}}
% \Pic is the operator that assigns to X its picard group, usage \Pic(X)
% \Picardstack_{X/B} denotes the Picard stack of X over B
% \Picardfunctor_{X/B} denotes the Picard functor of X over B
\def\Pic{\mathop{\mathrm{Pic}}\nolimits}
\def\Picardstack{\mathcal{P}\!\mathit{ic}}
\def\Picardfunctor{\mathrm{Pic}}
\def\Deformationcategory{\mathcal{D}\!\mathit{ef}}
\usepackage{cleveref}
\crefname{originalreference}{original reference}{original references}
\Crefname{originalreference}{Original reference}{Original references}
\makeatletter
\newlabel{lemma-find-extension}{{0GY2}{source}{}{sourcepointer-1}{}}
\newlabel{lemma-find-extension@cref}{{[originalreference][0][]0GY2}{[1][0][]source}{}{}{}}
\makeatother
\title{Two ideals with trivial sum can be interpolated by a unit-valued group element with prescribed residues}
\author{The Stacks Project Authors}
\date{stacks-project:etale-cohomology.tex:local-snapshot-20261011; GNU FDL 1.2 or later}
\begin{document}
\maketitle
\noindent\textit{Editorial note.} All delivered mathematical text is the author's own, copied from the pinned Stacks Project chapter file etale-cohomology.tex; no statement or proof step was rewritten, reordered or removed.\par\medskip
\noindent\textit{Editorial note.} The counted proof is the author's complete argument: the proof is written out in full: the choice of the element x and the change of ideals, the reformulation as a unit condition, and the application of the monogenic-extension lemma.\par\medskip
\noindent\textit{Editorial note.} The label lemma-interpolate names the counted statement, and the proof environment immediately following it in the source is delivered verbatim as its proof.\par\medskip
\noindent\textit{Editorial note.} The counted statement is carried by the author's own lemma environment in the chapter source.\par\medskip
\noindent\textit{Editorial note.} The statement and proof are copied from the pinned chapter file; the single cited input is kept as an external pointer with its canonical Stacks tag.\par\medskip
\noindent\textit{Editorial note.} The preamble is the chapter's own native preamble excerpt already used for this source: the Xy-pic setup (native preamble lines 14-18) and the native theorem declarations with the mathematical macros (native preamble lines 154-217).\par\medskip
\noindent\textit{Allowed standard prerequisites (editorial).}\begin{itemize}
\item Absolutely integrally closed A-algebras and the group G(R). Same chapter, the definition of an absolutely integrally closed ring and of the group G(R) of unit-valued elements used in the statement. Supplies the ambient ring, the ideals, the elements sigma\_0, sigma\_1 and the group in which g is sought.
\item The extension lemma for a pair of coprime ideals. Same chapter, the lemma on finding a monogenic extension with a unit of the form a - by, applied here to the reduction problem. Supplies the existence of the interpolating element once the problem is reduced to a unit condition.
\end{itemize}\medskip
\ifdefined\hypertarget\hypertarget{sourcepointer-1}{}\fi
\noindent\textit{Source reference.} Original reference 0GY2: sources/stacks-etale-cohomology/pinned-original-source.tex, source lines 14523--14523. Prerequisite pointer invoked at proof line(s) 14726; it names a standard fact of the same chapter that the counted argument uses, and every step of the counted proof built on it is retained verbatim in the delivered body.\par\medskip
\begin{lemma}
\label{lemma-interpolate}
Let $R$ be an $A$-algebra which is absolutely integrally closed.
Let $I, J \subset R$ be ideals with $I + J = R$.
There exists a $g \in G(R)$ such
that $g \bmod I = \sigma_0$ and $g \bmod J = \sigma_1$.
\end{lemma}

\begin{proof}
Choose $x \in I$ such that $x \equiv 1 \bmod J$.
We may and do replace $I$ by $xR$ and $J$ by $(x - 1)R$.
Then we are looking for an $s \in R$ such that
\begin{enumerate}
\item $1 + \pi s$ is a unit,
\item $s \equiv 0 \bmod xR$, and
\item $s \equiv 1 \bmod (x - 1)R$.
\end{enumerate}
The last two conditions say that $s = x + x(x - 1)y$ for some $y \in R$.
The first condition says that $1 + \pi s = 1 + \pi x + \pi x (x - 1) y$
needs to be a unit of $R$.
However, note that $1 + \pi x$ and $\pi x (x - 1)$ generate the
unit ideal of $R$ and that $1 + \pi x$ is an $\ell$th root of $1$
modulo $\pi x (x - 1)$\footnote{Because $1 + \pi x$ is congruent
to $1$ modulo $\pi$, congruent to $1$ modulo $x$, and congruent to
$1 + \pi = \zeta$ modulo $x - 1$ and because we have
$(\pi) \cap (x) \cap (x - 1) = (\pi x (x - 1))$ in $A[x]$.}.
Thus we win by Lemma \ref{lemma-find-extension} and the
fact that $R$ is absolutely integrally closed.
\end{proof}

\end{document}
